\begin{frame}[t]{\ev{\tagM\,\tagO}Three receipts misrepresented their events}
\vspace{0pt}
{\small
\begin{tabularx}{\textwidth}{@{}L{3.4cm}Y@{}}
\toprule
Recorded claim & What happened\\
\midrule
CI success & evaluation steps were skipped because a key was absent, and the job passed in 3\,s\\
Human approval & the harness answered the gate, and the record says \texttt{"mode": "human"}, actor \texttt{admin}\\
Structured CLI result & status text contaminated stdout, and an earlier attempt failed while parsing it\\
\bottomrule
\end{tabularx}\par}
\vspace{.35cm}
\begin{columns}[T,onlytextwidth]
\column{.48\textwidth}
\textbf{Required states}\\[2pt]
{\small executed\quad skipped\quad failed\par}
\column{.48\textwidth}
\textbf{Required actor}\\[2pt]
{\small human\quad harness\quad service\par}
\end{columns}
\vspace{.2cm}
{\small A trusted writer can still create a semantically incorrect receipt.\par}
\claim{A receipt must state what ran and who acted.}
\sources{Our CI, 28 Sep 2026; our work orders, Sep 2026: approval records and setup logs. Provenance: PASS (USENIX ATC'06); CamFlow (SoCC'17); in-toto (USENIX Security'19).}
\note{[0:40] Three of our own receipts misrepresented their events. A CI job reported success with its evaluation steps skipped. An approval record says human, but the harness answered. A command-line result mixed status text into its output. Writing receipts outside the sandbox protects them from the worker. Explicit states and actors make them interpretable. Provenance systems address the first part. Next we ask whether the record can reconstruct the run.}
\end{frame}
